% ======================================================================
% APPENDIX D: CORE ALGORITHM SOURCE CODE LISTINGS
% ======================================================================
\chapter{Core Algorithm Source Code Listings}
\label{app:listings}

\section{Dueling Q-Network and Agent Engine (\texttt{agent/dqn\_router.py})}
Listing~\ref{lst:dqn_router_code} presents the core PyTorch neural network class implementing Layer Normalization and dueling stream aggregation.

\begin{lstlisting}[language=Python, caption={PyTorch Dueling Q-Network Architecture with Layer Normalization}, label={lst:dqn_router_code}]
import torch
import torch.nn as nn
import torch.optim as optim
import numpy as np

class DuelingQNetwork(nn.Module):
    def __init__(self, state_size=10, action_size=4):
        super(DuelingQNetwork, self).__init__()
        # Shared Feature Extraction Trunk
        self.shared_fc1 = nn.Linear(state_size, 64)
        self.ln1 = nn.LayerNorm(64)
        self.shared_fc2 = nn.Linear(64, 128)
        self.ln2 = nn.LayerNorm(128)
        self.relu = nn.ReLU()

        # State-Value Stream V(s)
        self.val_fc1 = nn.Linear(128, 64)
        self.val_ln = nn.LayerNorm(64)
        self.val_out = nn.Linear(64, 1)

        # Action-Advantage Stream A(s, a)
        self.adv_fc1 = nn.Linear(128, 64)
        self.adv_ln = nn.LayerNorm(64)
        self.adv_out = nn.Linear(64, action_size)

    def forward(self, state):
        x = self.relu(self.ln1(self.shared_fc1(state)))
        x = self.relu(self.ln2(self.shared_fc2(x)))

        val = self.relu(self.val_ln(self.val_fc1(x)))
        val = self.val_out(val)

        adv = self.relu(self.adv_ln(self.adv_fc1(x)))
        adv = self.adv_out(adv)

        # Identifiability Aggregation: subtract mean advantage
        q_values = val + (adv - adv.mean(dim=-1, keepdim=True))
        return q_values
\end{lstlisting}

\section{Binary SumTree Prioritized Replay Buffer (\texttt{agent/prioritized\_replay.py})}
Listing~\ref{lst:sumtree_code} presents the complete non-recursive binary SumTree implementation.

\begin{lstlisting}[language=Python, caption={Iterative Binary SumTree and Prioritized Experience Replay}, label={lst:sumtree_code}]
import numpy as np

class SumTree:
    def __init__(self, capacity):
        self.capacity = capacity
        self.tree = np.zeros(2 * capacity - 1, dtype=np.float32)
        self.data = np.zeros(capacity, dtype=object)
        self.write_ptr = 0
        self.n_entries = 0

    def total(self):
        return self.tree[0]

    def add(self, p, data):
        idx = self.write_ptr + self.capacity - 1
        self.data[self.write_ptr] = data
        self.update(idx, p)
        self.write_ptr = (self.write_ptr + 1) % self.capacity
        if self.n_entries < self.capacity:
            self.n_entries += 1

    def update(self, idx, p):
        change = p - self.tree[idx]
        self.tree[idx] = p
        while idx > 0:
            idx = (idx - 1) // 2
            self.tree[idx] += change

    def get(self, s):
        idx = 0
        while idx < self.capacity - 1:
            left = 2 * idx + 1
            right = left + 1
            if s <= self.tree[left]:
                idx = left
            else:
                s -= self.tree[left]
                idx = right
        data_idx = idx - self.capacity + 1
        return idx, self.tree[idx], self.data[data_idx]
\end{lstlisting}

\section{State Feature Vector Extraction (\texttt{controller/state\_manager.py})}
Listing~\ref{lst:state_features_code} shows the extraction of the 10-dimensional state vector from the active NetworkX graph.

\begin{lstlisting}[language=Python, caption={Normalized 10-Dimensional State Extraction in StateManager}, label={lst:state_features_code}]
def get_routing_state_features(self, candidate_paths):
    state = np.zeros(10, dtype=np.float32)
    if not candidate_paths or len(candidate_paths) == 0:
        return state

    # Feature 0: Normalized hop count of shortest path
    shortest_path = candidate_paths[0]
    hops = max(1, len(shortest_path) - 1)
    state[0] = np.clip(hops / 10.0, 0.1, 1.0)

    # Features 1-4: Peak link utilization on up to 4 candidate paths
    # Features 5-8: Normalized cumulative propagation latency
    for i in range(min(4, len(candidate_paths))):
        path = candidate_paths[i]
        path_util = 0.0
        path_lat = 0.0
        for u, v in zip(path[:-1], path[1:]):
            if self.graph.has_edge(u, v):
                edge_data = self.graph[u][v]
                path_util = max(path_util, edge_data.get('utilization', 0.0))
                path_lat += edge_data.get('latency', 2.0)
        state[1 + i] = np.clip(path_util, 0.0, 1.0)
        state[5 + i] = np.clip(path_lat / 50.0, 0.0, 1.0)

    # Feature 9: Global peak link utilization across entire network
    max_util = 0.0
    for u, v, data in self.graph.edges(data=True):
        max_util = max(max_util, data.get('utilization', 0.0))
    state[9] = np.clip(max_util, 0.0, 1.0)

    return state
\end{lstlisting}
